\section{The published Maynard core}
\label{sec:published-core}

We first specialize the abstract assembly theorem using only the
formally published distribution input.  Write
\begin{equation}
 \lambda=\log_XL,
 \qquad \nu=\log_XD_0.
 \label{eq:maynard-exponent-coordinates}
\end{equation}
Ignoring fixed power margins, the Maynard region is the polytope
\begin{equation}
 2\lambda+\nu\le1,
 \qquad
 7\lambda+12\nu\le4,
 \qquad
 19\lambda+20\nu\le10.
 \label{eq:maynard-exponent-polytope}
\end{equation}
The first and third faces meet at
\begin{equation}
 (\lambda,\nu)=\left(\frac{10}{21},\frac1{21}\right),
 \qquad \lambda+\nu=\frac{11}{21},
 \label{eq:maynard-assembly-vertex}
\end{equation}
and the second inequality is then strict.

\begin{theorem}[Published assembled wedge]
\label{thm:published-assembled-wedge}
Fix $h\ne0$, $0<\delta<1/2$, and $\eta>0$.  Suppose
\eqref{eq:assembly-geometry} and the finite-row leakage condition in
\eqref{eq:assembly-analysis} hold, together with
\begin{equation}
 D_0L^2\le X^{1-\eta},
 \qquad
 D_0^{12}L^7\le X^{4-\eta},
 \qquad
 D_0^{20}L^{19}\le X^{10-\eta}.
 \label{eq:published-assembly-conditions}
\end{equation}
Then, for every $A>0$,
\begin{equation}
 \cC_{L,K}(D_0)\ll_A X(\log X)^{-A},
 \qquad
 \cH^{\mathrm{HB}}_{L,K}
 =\cT_{L,K}(D_0,V)+O_A(X(\log X)^{-A}).
 \label{eq:published-assembled-conclusion}
\end{equation}
All implied constants may depend on the fixed data listed in
\cref{thm:assembled-prefix}.
\end{theorem}

\begin{proof}
The three inequalities place $(L,D_0)$ in
$\mathfrak R^{\mathrm M}_\eta$.  Apply
\cref{thm:assembled-prefix,cor:tail-normal-form}.
\end{proof}

The following explicit family reaches the vertex from the interior.

\begin{corollary}[A complete prefix near the
\texorpdfstring{$11/21$}{11/21} boundary]
\label{cor:published-explicit-prefix}
Fix $0<\sigma<1/100$ and
\begin{equation}
 \frac1{42}+3\sigma<\delta<\frac14.
 \label{eq:published-delta-range}
\end{equation}
Set, up to harmless integer parts,
\begin{equation}
 L=X^{10/21-\sigma},
 \qquad
 D_0=X^{1/21-\sigma},
 \qquad
 K=X/L=X^{11/21+\sigma}.
 \label{eq:published-explicit-scales}
\end{equation}
Then \eqref{eq:published-assembled-conclusion} holds for every $A>0$.
Thus all divisor labels
\begin{equation}
 d\le X^{1/21-\sigma}
 \label{eq:published-prefix-length}
\end{equation}
are removed in one cumulative packet.
\end{corollary}

\begin{proof}
The three monomials in
\eqref{eq:published-assembly-conditions} have exponents
\begin{equation}
 1-3\sigma,
 \qquad
 \frac{82}{21}-19\sigma,
 \qquad
 10-39\sigma.
 \label{eq:published-monomial-check}
\end{equation}
They therefore have fixed power margins.  Moreover,
\begin{equation}
 LD_0R
 =X^{\,1+1/42-2\sigma-\delta+o(1)}
 \le X^{1-5\sigma+o(1)}
 \label{eq:published-leakage-check}
\end{equation}
by \eqref{eq:published-delta-range}.  The same lower bound on $\delta$
gives $L>2R$ for large $X$.  The upper bound $\delta<1/4$ gives
$D_0<V$, while $K>2V$ is immediate.  Hence all hypotheses of
\cref{thm:published-assembled-wedge} hold after choosing a fixed
$\eta>0$ smaller than the displayed margins.
\end{proof}

\begin{remark}[The cumulative packaging of TPC-16 slices]
TPC-16 proved the corresponding estimate for one smooth divisor slice
$d\asymp D$.  Under uniform fixed margins, one could also sum
$O(\log X)$ such estimates after requesting additional logarithmic
saving.  The statement here instead inserts
$\mu(d)\one_{d\le D_0}$ itself in one application and records the exact
tail normal form.  It does not increase the published exponent vertex.
\end{remark}
